\section{模型评价与改进}
\subsection{模型的优点}
\begin{enumerate}[label=\arabic*.]
\item 有限体积法使公共界面两侧通量大小相等、方向相反，内部通量在全域求和时相互抵消，便于检查离散方程的收支一致性。
\item 四个问题依次加入温湿耦合、全域终止事件和动态几何更新，共用网格与边界处理方式，计算框架具有较强的复用性。
\item 二维轴对称模型同时描述侧壁与端面的热量和水分交换，并通过全域扫描确定最湿位置。
\item 近壁加密、网格收敛性比较及环境延拓敏感性分析，分别检验了高梯度区域离散、网格规模和长时边界选择对结果的影响。

\end{enumerate}

\subsection{模型的不足与改进方向}
\begin{enumerate}[label=\arabic*.]
\item 温度方程未计入水分蒸发的耗热效应，后续可增加潜热项及相应的热质耦合关系。
\item 模型未反映药材纤维与孔隙结构的方向差异，后续可引入方向相关的导热系数和水分扩散系数。
\item 数据截止后的环境由末段均值给定，所考察方案间的差异较小不能替代长时间实测边界数据。
\item 问题四按给定半径轨迹更新几何，后续应完善干物质密度与骨架运动的相容关系，并在相同参数关系下增加固定半径对照，以单独量化收缩影响。
\end{enumerate}
